\section{Introduction}
\label{sec:introduction}
% 第一章：引言

% 第1段：
% 介绍研究背景和实际应用场景，阐述所研究的调度或决策问题
% 在智能制造、柔性生产等实际应用中的重要性和现实需求。


% 第2段：
% 描述该问题面临的核心挑战，例如：
% 生产过程柔性、多阶段加工与装配耦合、动态扰动、
% 不确定性因素、多重约束以及实时在线决策需求。


% 第3段：
% 总结现有传统优化方法以及基于向量状态表示的深度强化学习方法的不足。
% 重点分析其在复杂制造关系建模、动态环境适应、
% 问题规模扩展以及实时响应能力方面存在的局限。


% 第4段：
% 说明图结构表示方法和PPO强化学习方法适用于该问题的原因。
% 进一步指出现有研究仍存在的具体不足，
% 明确本文需要解决的研究空白（research gap）。


% 第5段：
% 简要介绍本文提出的\method{}框架及整体设计思想。
% 说明该方法如何通过状态表示、特征编码和策略优化实现动态调度决策。

The main contributions should be stated as verifiable claims:
\begin{itemize}
  \item {[Contribution 1: problem formulation or decision framework.]}
  \item {[Contribution 2: heterogeneous graph construction or encoder design.]}
  \item {[Contribution 3: PPO policy, reward, action, or training design.]}
  \item {[Contribution 4: experimental evidence and practical significance.]}
\end{itemize}

% 文章最后给出全文组织结构说明。
% 各章节安排应与本文实际内容保持一致，
% 简要介绍后续章节分别介绍的问题建模、方法设计、实验验证和结论分析。
